\section{Grundlagen}\label{sec:Grundlagen}
Hier werden die Grundlagen zur Hairpin Technologie und der Automatisierung beschrieben.


\subsection{Hairpin-Technologie}\label{sec:HairPinTechnologie}

\subsubsection{Entwicklung}\label{sec:EntwicklungHairPin}
Wie und waraus hat sich die Hairpin-Technologie entwickelt? Welche Vorteile bietet sie gegenüber anderen Technologien? Welche Nachteile gibt es?

\subsubsection{Produktionsablauf}\label{sec:ProduktionsablaufHairPin}
Wie sieht der Produktionsablauf aus? Welche unterschiede gibt es zu Prototypen Fertigung und Serienfertigung? Welcher Abschnitt wird in dieser Arbeit betrachtet?


\subsection{Produktionsverfahren zum Abisolieren}\label{sec:ProduktionsverfahrenAbisolieren}
Welche Verfahren gitb es? Wo sind diese in der Produkitonskette angesiedelt?

\subsubsection{Fräsen Inline}\label{sec:FräsenInline}
Serienproduktion, Vorteile, Nachteile, Kosten, Taktzeit, Qualität, Anlage im PEM?

\subsubsection{Laser Inline}\label{sec:LaserInline}
Serienproduktion, Vorteile, Nachteile, Kosten, Taktzeit, Qualität, Anlage im PEM?

\subsubsection{Laser als eigener Prozess}\label{sec:LaserEigener}
Serienproduktion, Vorteile, Nachteile, Kosten, Taktzeit, Qualität, etc.


\subsection{Normen und Vorgaben}\label{sec:NormenVorgaben}
Welche Normen muss ich beachten? Was gibts zu Sicherheit beim Anlagenbau?